\chapter{Hedge funds quantitatifs et gérants systématiques}\label{ch:in:quantitative-hedge-funds-and-systematic-managers}

Les fonds les plus discrets du métier déposent malgré tout un formulaire public. Chaque conseiller en investissement enregistré auprès de la SEC
répond, au moins une fois par an et sous peine de sanctions légales, aux mêmes questions~: combien de personnes il emploie, combien de salariés
exercent une activité d'investissement, combien il gère et combien de fonds privés il administre, avec leurs actifs bruts. La SEC
publie les réponses dans un fichier unique, une ligne par conseiller. En janvier 2026, ce fichier recensait 2\,577 conseillers qui gèrent des
hedge funds. Lues ensemble, leurs lignes rendent comparables des sociétés qui ne publient rien d'autre~: le conseiller médian emploie onze
personnes, et parmi ceux qui comptent dix salariés ou plus, la moitié centrale gère entre 58~millions et 221~millions de dollars par
salarié. Ce chapitre examine les gérants quantitatifs parmi eux.

\section{Trois familles}

Les gérants quantitatifs sont généralement répartis en trois familles, qui diffèrent par la méthode, l'horizon et le
type d'emploi de recherche qu'ils offrent.
\begin{itemize}
\item \textbf{Fonds multistratégies pilotés par la recherche} exploitent de nombreux signaux sur plusieurs classes d'actifs et horizons, souvent
aux côtés d'équipes discrétionnaires, sur de grandes plateformes de recherche.
\item \textbf{Suiveurs de tendance}, les commodity trading advisors du livre~1 (chapitre~3), négocient des contrats à terme sur la persistance
des mouvements de prix (livre~8, chapitre~19), sur des horizons allant de quelques semaines à quelques mois.
\item \textbf{Maisons d'arbitrage statistique} négocient de vastes portefeuilles d'actions neutres au marché, sur des signaux relatifs de court
horizon (livre~8, chapitre~1).
\end{itemize}
Ces familles se recoupent, et les sociétés n'emploient pas toujours ces termes pour elles-mêmes. Ce qui distingue les gérants de ce
chapitre des autres, c'est la manière dont les décisions sont prises.

\begin{definition}[Gérant systématique, gérant discrétionnaire]\label{def:in:quantitative-hedge-funds-and-systematic-managers:sys}
Un \emph{gérant systématique}\index{gérant systématique} est un gérant d'investissement dont les décisions de trading sont prises par des règles
implémentées dans un logiciel, à partir de modèles calibrés sur des données~; des personnes conçoivent, testent, surveillent et modifient les règles, et non les
opérations. Un \emph{gérant discrétionnaire}\index{gérant discrétionnaire} est un gérant dont les gérants de portefeuille décident des opérations,
quelle que soit l'analyse qui les étaye.
\end{definition}

Les gérants cités dans ce chapitre se décrivent eux-mêmes en ces termes~: une entreprise d'investissement qui «~emploie des méthodes mathématiques
et statistiques dans la conception et l'exécution de ses programmes d'investissement~», «~une société quantitative d'investissement et de trading~»,
«~investisseurs quantitatifs~» pour qui «~une approche systématique et disciplinée est la meilleure voie~», «~l'un des
plus anciens \omterm{def:in:quantitative-hedge-funds-and-systematic-managers:sys}{gérants systématiques}~», une société qui aborde l'investissement «~à travers le prisme de l'apprentissage automatique~».
L'une d'elles déclare avoir «~été lancée en 1988 comme pionnière de l'investissement systématique~», puis avoir développé des compétences fondamentales et discrétionnaires,
car de nombreuses grandes sociétés cumulent les deux.

\section{Ce que révèle le Form ADV}

\begin{definition}[Encours réglementaires sous gestion]\label{def:in:quantitative-hedge-funds-and-systematic-managers:raum}
Les \emph{encours réglementaires sous gestion}\index{encours réglementaires sous gestion} (RAUM) d'un conseiller sont les portefeuilles de titres
sur lesquels il fournit des services continus et réguliers de surveillance ou de gestion, à leur valeur brute~: le
conseiller ne déduit aucun endettement en cours ni aucun autre passif des portefeuilles.
\end{definition}

Les RAUM ne sont pas les encours sous gestion du livre~1 (chapitre~1) tels que les fonds les déclarent à leurs investisseurs. Ils comptent les positions brutes
donc un fonds qui emprunte ou détient des dérivés déclare plus que sa valeur liquidative~; et ils ne comptent que
des portefeuilles de titres. Un gérant au capital modeste mais à fort effet de levier brut a des RAUM élevés. Pour un candidat à l'emploi, cet écart
compte lorsqu'on compare des sociétés~: les RAUM mesurent ce que le conseiller gère, et non ce que les investisseurs lui ont confié.

\begin{definition}[Conseiller exempté déclarant]\label{def:in:quantitative-hedge-funds-and-systematic-managers:era}
Un \emph{conseiller exempté déclarant}\index{conseiller exempté déclarant} (exempt reporting adviser) est un conseiller en investissement dispensé d'enregistrement auprès de la SEC --
par exemple parce qu'il ne conseille que des fonds privés et gère moins de 150~millions de dollars d'actifs de fonds privés aux
États-Unis -- qui doit malgré tout déposer un bref rapport sur le Form ADV.
\end{definition}

Quatre rubriques du formulaire contiennent l'essentiel de ce qu'un lecteur cherche~:
\begin{itemize}
\item \textbf{5.A} le nombre approximatif de salariés~;
\item \textbf{5.B(1)} combien d'entre eux «~exercent des fonctions de conseil en investissement (y compris la recherche)~»~;
\item \textbf{5.F} \omterm{def:in:quantitative-hedge-funds-and-systematic-managers:raum}{encours réglementaires sous gestion}~;
\item \textbf{7.B} les fonds privés conseillés~: si l'un d'entre eux est un hedge fund, combien ils sont et le total de leurs actifs bruts.
\end{itemize}

\begin{remark}[Ce que le formulaire ne montre pas]\label{rem:in:quantitative-hedge-funds-and-systematic-managers:limits}
Le formulaire décrit un seul conseiller enregistré. Un groupe peut répartir son personnel et ses actifs entre plusieurs conseillers et
entités affiliées, dont certaines enregistrées ailleurs~: le conseiller américain d'un groupe mondial n'emploie peut-être que quelques dizaines de ses milliers de
salariés. Le formulaire ne publie ni rendements, ni frais, ni rémunérations, ni chiffre d'affaires. Les réponses sont mises à jour
au moins une fois par an, de sorte qu'une ligne peut dater de presque un an.
\end{remark}

\section{Actifs, actifs bruts et concentration}

Parmi les 1\,458 conseillers de hedge funds comptant dix salariés ou plus dans le fichier de janvier 2026, les RAUM par salarié ont une
médiane de 119~millions de dollars, une moitié centrale comprise entre 58~millions et 221~millions de dollars, et un dixième des conseillers au-dessus de 399~millions de dollars
(voir la \cref{fig:in:quantitative-hedge-funds-and-systematic-managers:hist}). La moitié des salariés d'un conseiller type exerce une activité d'investissement
(les quartiles de la part de l'effectif consacré au conseil sont de 37~\% et 62~\%). Le secteur n'est pas concentré~: les dix plus grands
conseillers selon leurs actifs bruts de fonds privés détiennent 21~\% des 17.6~billions de dollars déclarés par l'ensemble des conseillers de hedge funds, et
l'indice de Herfindahl--Hirschman (livre~16, chapitre~30) de ces actifs bruts vaut 78 sur une échelle de 10\,000.

\begin{omfigure}
\begin{tikzpicture}
\begin{axis}[width=11cm,height=5.4cm,ybar,bar width=7pt,xmin=5.8,xmax=10.3,ymin=0,ymax=380,
  xlabel={\small RAUM par salarié (échelle logarithmique, \$)},ylabel={\small conseillers},
  xtick={6,7,8,9,10},xticklabels={1 million,10 millions,100 millions,1 milliard,10 milliards},
  tick label style={font=\footnotesize},grid=major,grid style={gray!20},table/col sep=comma]
\addplot[fill=omDef!55,draw=omDef] table[x=mid,y=advisers]{figdata/industry/04-quantitative-hedge-funds-and-systematic-managers/rpe_hist.csv};
\draw[omThm,thick,dashed] (axis cs:7.766,0) -- (axis cs:7.766,360);
\draw[omThm,thick] (axis cs:8.077,0) -- (axis cs:8.077,360);
\draw[omThm,thick,dashed] (axis cs:8.345,0) -- (axis cs:8.345,360);
\end{axis}
\end{tikzpicture}

\omcaption{\omterm{def:in:quantitative-hedge-funds-and-systematic-managers:raum}{Encours réglementaires sous gestion} par salarié des 1\,458 conseillers de hedge funds enregistrés auprès de la SEC comptant dix
ou plus, en janvier 2026, sur une échelle logarithmique (classes d'un quart de décade). Trait plein~: médiane, 119~millions de dollars~; pointillés~:
quartiles, 58~millions et 221~millions de dollars. Données~: fichier des conseillers de la SEC, rubriques 5.A et 5.F, via
\texttt{in\_quantfunds.hist}.}\label{fig:in:quantitative-hedge-funds-and-systematic-managers:hist}
\end{omfigure}

\begin{dated}{2026-01}{Treize conseillers quantitatifs dans le fichier de la SEC}\label{dat:in:quantitative-hedge-funds-and-systematic-managers:named}
\small
\begin{tabular}{@{}>{\raggedright\arraybackslash}p{3.8cm}>{\centering\arraybackslash}p{1.7cm}>{\centering\arraybackslash}p{1.7cm}>{\centering\arraybackslash}p{2.0cm}>{\centering\arraybackslash}p{3.0cm}@{}}
\toprule
conseiller (tel qu'enregistré) & salariés & dont conseil & RAUM (Md\$) & actifs bruts des fonds privés (Md\$)\\
\midrule
Squarepoint (conseiller américain) & 71 & 51 & 182.2 & 162.5\\
AQR Capital Management & 611 & 238 & 160.5 & 87.6\\
D. E. Shaw \& Co. & 928 & 476 & 154.6 & 161.1\\
Two Sigma Investments & 1\,575 & 749 & 110.3 & 139.7\\
Renaissance Technologies & 300 & 160 & 92.0 & 81.3\\
Two Sigma Advisers & 101 & 68 & 74.3 & 107.8\\
AHL Partners (Man Group) & 162 & 131 & 70.7 & 56.5\\
Graham Capital Management & 204 & 98 & 28.5 & 28.5\\
Voleon Capital Management & 111 & 18 & 16.4 & 17.3\\
PDT Partners & 270 & 92 & 14.0 & 14.4\\
Capital Fund Management & 285 & 79 & 12.3 & 12.4\\
Winton Capital Management & 146 & 35 & 10.1 & 4.6\\
Aspect Capital & 142 & 60 & 3.5 & 3.4\\
\bottomrule
\end{tabular}

\smallskip\noindent Les derniers dépôts Form ADV entre mars et décembre 2025, tels que publiés dans le fichier de la SEC de janvier
2026. Chaque ligne correspond à un conseiller enregistré et non à l'entreprise~: les groupes déposent par l'intermédiaire de plusieurs
conseillers et emploient du personnel en dehors d'eux.
\end{dated}

Le tableau met côte à côte des entreprises de formes très différentes
(la \cref{fig:in:quantitative-hedge-funds-and-systematic-managers:scatter}). Un conseiller déclare 182~milliards de dollars de RAUM avec
71 salariés, soit 2.6~milliards de dollars chacun~: ses équipes travaillent pour des filiales hors de l'entité enregistrée, si bien que
le ratio mesure l'entité et non l'entreprise. Deux conseillers d'un même groupe déposent séparément, avec 1\,575 et 101 salariés. Les additionner compterait
deux fois tout portefeuille qu'un des deux sous-conseille pour l'autre. À l'autre extrémité, un gérant de suivi de tendance déclare 24~millions de dollars par
salarié, et un gérant d'apprentissage automatique compte 18 de ses 111 salariés dans des fonctions de conseil~: les autres conçoivent et font
tourner les systèmes. La part des fonctions de conseil est la réponse la plus proche du formulaire à la question qu'un candidat se pose~: combien de membres du personnel font de la recherche.

\begin{omfigure}
\begin{tikzpicture}
\begin{axis}[width=10.5cm,height=6.4cm,xmode=log,ymode=log,xmin=40,xmax=3000,ymin=2,ymax=400,
  xlabel={\small salariés (échelle logarithmique)},ylabel={\small RAUM, milliards de \$ (échelle logarithmique)},
  tick label style={font=\footnotesize},grid=major,grid style={gray!20},table/col sep=comma,
  log ticks with fixed point]
\addplot[only marks,mark=*,mark size=2pt,omDef] table[x=employees,y=raum_bn]{figdata/industry/04-quantitative-hedge-funds-and-systematic-managers/named.csv};
\addplot[omThm,dashed,domain=40:3000,samples=2] {0.119*x};
\addplot[gray,dotted,domain=40:3000,samples=2] {0.0583*x};
\addplot[gray,dotted,domain=40:3000,samples=2] {0.221*x};
\node[font=\scriptsize,anchor=west] at (axis cs:71,182.2) {Conseiller américain d'un groupe mondial};
\end{axis}
\end{tikzpicture}

\omcaption{Les treize conseillers de l'encadré daté~: salariés contre RAUM, tous deux en échelles logarithmiques. La ligne tiretée est la
médiane du RAUM par salarié de tous les conseillers de hedge funds comptant dix salariés ou plus, les pointillés en donnent les quartiles.
Données~: fichier des conseillers de la SEC de janvier 2026, via \texttt{in\_quantfunds.named}.}\label{fig:in:quantitative-hedge-funds-and-systematic-managers:scatter}
\end{omfigure}

\begin{method}[Lire le formulaire ADV d'un gérant]\label{met:in:quantitative-hedge-funds-and-systematic-managers:read}
\begin{enumerate}
\item Repérer chaque conseiller du groupe (les dénominations légales du fichier, les déclarations propres de l'entreprise) et lire leurs
lignes séparément.
\item Comparer le RAUM aux actifs bruts des fonds privés~: un ratio très inférieur à un signifie des comptes gérés ou des fonds enregistrés
hors des fonds privés~; très supérieur à un, des fonds privés comptés dans le RAUM d'un autre conseiller.
\item Lire la part des fonctions de conseil (5.B(1) sur 5.A) comme la part du personnel affecté à la recherche et au travail de portefeuille.
\item Vérifier la date de dépôt~; traiter chaque chiffre comme un intervalle daté.
\end{enumerate}
\end{method}

\section{Secret et publication}

Les \omterm{def:in:quantitative-hedge-funds-and-systematic-managers:sys}{gérants systématiques} sont discrets sur leurs modèles et ouverts sur leurs méthodes. Ce qu'ils publient relève de trois
catégories~: des articles de recherche et des commentaires sur les facteurs, les tendances et les primes de risque, écrits pour les investisseurs~; des articles académiques signés par
leurs chercheurs, souvent en statistique et en apprentissage automatique~; et des descriptions de méthode sur leurs propres pages, comme
la déclaration d'une entreprise selon laquelle elle utilise «~des méthodes scientifiques et académiques, notamment les dernières techniques d'apprentissage automatique
pour détecter des signaux de trading dans les données financières mondiales et alternatives~». Rien de tout cela ne révèle un signal. Cela renseigne le candidat sur
ce que l'entreprise valorise et sur le type de poste de recherche qu'elle propose~: un gérant qui publie des travaux de recherche sur les facteurs recrute des chercheurs qui
écrivent~; un gérant qui ne publie qu'une phrase recrute des personnes qui n'écriront rien.

\section{Tutoriel~: RAUM par salarié à partir du fichier des conseillers de la SEC}\label{tut:in:quantitative-hedge-funds-and-systematic-managers:tut}

\begin{tutorial}
\textbf{Objectif.} Calculer la distribution du RAUM par salarié des conseillers de hedge funds et situer treize conseillers quantitatifs
nommés dans cette distribution. \textbf{État final~:} la \cref{fig:in:quantitative-hedge-funds-and-systematic-managers:hist} et les rangs de
l'exercice~5.
\begin{tutsteps}
\item \textbf{Le fichier.} Télécharger le fichier des conseillers de la SEC pour un mois (environ 5\,Mo compressés, un CSV d'environ 16\,600 lignes
et 448 colonnes) dans un répertoire temporaire~; il est public et n'est pas versionné.
\item \textbf{Le lire.} \texttt{firm.formadv.read(path)} parcourt le CSV depuis l'archive zip et conserve, pour les conseillers de hedge
funds, les rubriques de la section~2 (le \cref{lst:in:quantitative-hedge-funds-and-systematic-managers:read}).
\omcode{code/firm/formadv/firm_formadv.py}{81}{97}{Parcourir le fichier des conseillers et conserver ceux des hedge funds.}\label{lst:in:quantitative-hedge-funds-and-systematic-managers:read}
\item \textbf{Dériver.} \texttt{in\_adv\_derive.py} écrit trois petits tableaux dans \texttt{data/industry/}~: les
statistiques de distribution, les lignes des conseillers nommés et un histogramme~; les figures et les tests du chapitre ne lisent que ceux-ci.
\item \textbf{Situer.} \texttt{in\_quantfunds.percentile\_rank(x)} donne la part des conseillers comptant dix salariés ou plus
dont le RAUM par salarié est inférieur à $x$.
\end{tutsteps}
\end{tutorial}

\textbf{Pour aller plus loin.} Refaire le calcul avec le fichier d'un an plus tôt et comparer les salariés des conseillers nommés
(le chapitre~5 le fait pour les plateformes multigérants)~; restreindre aux conseillers dont les actifs bruts des fonds privés dépassent
leur RAUM et observer comment la distribution se déplace.

\section{Construction~: le lecteur du fichier des conseillers}\label{bld:in:quantitative-hedge-funds-and-systematic-managers:formadv}

\begin{build}
\bfield{Rôle} Un seul lecteur pour le fichier Form ADV de la SEC, utilisé par les chapitres 4, 5, 8 et 12.
\bfield{Interface} \texttt{firm.formadv}~: \texttt{COLS}~;
\texttt{Adviser(crd, name, employees, advisory, raum, pf\_gav, n\_hedge, filed)}~; \texttt{read(path\_or\_fileobj, hedge\_only)}~; \texttt{per\_head(a)}~;
\texttt{quantiles(values, qs)}~; \texttt{concentration(values, k)}, qui enveloppe \texttt{firm.moats}.
\bfield{Règles} Parcourt le fichier et n'écrit jamais les lignes brutes~; une réponse vide vaut \texttt{None}, jamais zéro~; les
ratios par salarié ne sont pas définis sans un effectif positif.
\bfield{Tests d'acceptation} \texttt{code/firm/formadv/tests/}~: un jeu de trois lignes au format du fichier~; le filtre des hedge
funds~; les ratios par salarié et un nombre de conseillers manquant~; la part des $k$ premiers et l'indice sur des parts connues.
\bfield{Pour aller plus loin} Lire les fichiers mensuels de plusieurs années dans un panel indexé par le numéro CRD du conseiller~; lire les
lignes de fonds privés de l'annexe D à partir du jeu complet de données du Form ADV.
\end{build}

\begin{omsources}
\item SEC, Information about registered investment advisers, fichier des conseillers de janvier 2026 (\texttt{ia010226.zip}),
copie de l'Internet Archive.
\item SEC, Form ADV Part 1A and General Instructions; 17 CFR 275.203(m)-1 and 275.204-4.
\item Les pages propres des gérants pour leurs descriptions (ledger F7--F18).
\end{omsources}

\section{Exercices}

\begin{exercise}[$\star$]\label{exo:in:quantitative-hedge-funds-and-systematic-managers:1}
Calculer le RAUM par salarié et la part des fonctions de conseil du conseiller qui compte 300 salariés, 160 salariés affectés au conseil et
92.0~milliards de dollars de RAUM.
\end{exercise}

\begin{exercise}[$\star$]\label{exo:in:quantitative-hedge-funds-and-systematic-managers:2}
Un fonds a des actifs nets de 2~milliards de dollars et des positions brutes longues et courtes en titres de 7~milliards de dollars. Quel RAUM
son conseiller déclare-t-il à peu près pour lui, et pourquoi~?
\end{exercise}

\begin{exercise}[$\star$]\label{exo:in:quantitative-hedge-funds-and-systematic-managers:3}
Un conseiller ne gère que des fonds privés, avec 120~millions de dollars d'actifs aux États-Unis. Doit-il s'enregistrer auprès de la
SEC, et que doit-il déposer~?
\end{exercise}

\begin{exercise}[$\star\star$]\label{exo:in:quantitative-hedge-funds-and-systematic-managers:4}
Parmi les treize conseillers de l'encadré daté, lesquels ont des actifs bruts de fonds privés inférieurs à la moitié de leur RAUM, et que
suggère ce résultat~?
\end{exercise}

\begin{exercise}[$\star\star$]\label{exo:in:quantitative-hedge-funds-and-systematic-managers:5}
En utilisant \texttt{percentile\_rank}, où se situent dans la distribution du RAUM par salarié le plus petit et le plus grand des
treize conseillers selon cette mesure, et pourquoi le plus grand n'est-il pas un fait concernant sa société~?
\end{exercise}

\begin{exercise}[$\star\star$]\label{exo:in:quantitative-hedge-funds-and-systematic-managers:6}
À partir de la \cref{fig:in:quantitative-hedge-funds-and-systematic-managers:hist}, estimer la part des conseillers dont le RAUM par
salarié dépasse 1~milliard de dollars.
\end{exercise}

\begin{exercise}[$\star\star\star$]\label{exo:in:quantitative-hedge-funds-and-systematic-managers:7}
\emph{Programmation.} Calculer la part des dix premiers et l'indice de Herfindahl--Hirschman de dix conseillers de taille égale et de mille
conseillers de taille égale. Où se situe l'indice de 78 du fichier entre ces deux valeurs, et que signifie-t-il~?
\end{exercise}

\begin{exercise}[$\star\star\star$]\label{exo:in:quantitative-hedge-funds-and-systematic-managers:8}
\emph{Trouver la faille.} «~Le conseiller de 71 salariés gère 182~milliards de dollars~: c'est le plus productif des gérants quantitatifs
du monde.~»
\end{exercise}

\section{Problème~: actifs par tête}

\begin{problem}[{Problème du week-end --- actifs par tête}]\label{pb:in:quantitative-hedge-funds-and-systematic-managers:1}
Un chercheur qui pèse les offres de trois \omterm{def:in:quantitative-hedge-funds-and-systematic-managers:sys}{gérants systématiques} veut savoir la taille de chacun et la part de leur activité
tournée vers la recherche, à partir des seules données publiques.

\medskip\noindent\textbf{Partie I --- Les familles et le formulaire.}
\begin{enumerate}
\item Définir un \omterm{def:in:quantitative-hedge-funds-and-systematic-managers:sys}{gérant systématique} et un \omterm{def:in:quantitative-hedge-funds-and-systematic-managers:sys}{gérant discrétionnaire}, et décrire les trois familles.
\item Quels quatre éléments du formulaire ADV portent l'essentiel de ce qu'un lecteur cherche~?
\item Définir le RAUM et expliquer en quoi il diffère des encours sous gestion.
\item Définir un \omterm{def:in:quantitative-hedge-funds-and-systematic-managers:era}{conseiller exempté déclarant} et donner le seuil de l'exemption pour fonds privés.
\item Quelles trois limites du formulaire la \cref{rem:in:quantitative-hedge-funds-and-systematic-managers:limits} énonce-t-elle~?
\end{enumerate}

\medskip\noindent\textbf{Partie II --- La distribution.}
\begin{enumerate}[resume]
\item Combien de conseillers de hedge funds figurent dans le fichier de janvier 2026, et combien en ont dix salariés ou plus~?
\item Donner la médiane et les quartiles du RAUM par salarié pour ces derniers.
\item Donner les quartiles de la part de conseil.
\item Donner la part des dix plus grands conseillers dans les actifs bruts de fonds privés, ainsi que l'indice de Herfindahl--Hirschman.
\item Quelle est la médiane du nombre de salariés d'un conseiller de hedge fund~?
\end{enumerate}

\medskip\noindent\textbf{Partie III --- Les conseillers nommés.}
\begin{enumerate}[resume]
\item Calculer le RAUM par salarié des treize conseillers nommés et les classer.
\item Lesquels se situent dans la moitié centrale de la distribution~?
\item Lesquels ont la part de conseil la plus faible et la plus élevée~?
\item Quelles deux lignes appartiennent à un même groupe, et pourquoi ne faut-il pas les additionner~?
\item Quel rapport entre RAUM et actifs bruts de fonds privés est le plus éloigné de un~?
\end{enumerate}

\medskip\noindent\textbf{Partie IV --- Le verdict.}
\begin{enumerate}[resume]
\item Énoncer le \emph{résultat nommé}~: la moitié centrale du RAUM par salarié pour les conseillers de hedge funds, l'étendue des
conseillers quantitatifs nommés et les mesures de concentration.
\item Pourquoi le plus grand rapport nommé est-il un fait concernant une entité et non une société~?
\item Que dit une part de conseil de 16~\% à un chercheur sur une société~?
\item De quoi le chercheur aurait-il besoin que le formulaire ne publie pas~?
\item En deux phrases, comment le chercheur doit-il utiliser les trois lignes~?
\end{enumerate}
\end{problem}

\section{Questions d'entretien}

\begin{interviewq}[$\star$ \iqroles{chercheur}]\label{iq:in:quantitative-hedge-funds-and-systematic-managers:1}
Quelle est la différence entre un fonds systématique et un fonds discrétionnaire, du point de vue
d'un chercheur qui y travaille~?
\end{interviewq}

\begin{interviewq}[$\star$ \iqroles{chercheur, risque}]\label{iq:in:quantitative-hedge-funds-and-systematic-managers:2}
Pourquoi les actifs bruts d'un fonds peuvent-ils dépasser sa valeur liquidative de plusieurs fois~?
\end{interviewq}

\begin{interviewq}[$\star\star$ \iqroles{chercheur}]\label{iq:in:quantitative-hedge-funds-and-systematic-managers:3}
Un gérant affiche 100~milliards de dollars de RAUM et 1\,000 salariés, dont 400 font du travail d'investissement. Comment le décririez-vous, et
que voudriez-vous savoir avant de le comparer à une société de 300 personnes gérant 90~milliards de dollars~?
\end{interviewq}

\begin{interviewq}[$\star\star$ \iqroles{chercheur, ingénieur ML}]\label{iq:in:quantitative-hedge-funds-and-systematic-managers:4}
Pourquoi les \omterm{def:in:quantitative-hedge-funds-and-systematic-managers:sys}{gérants systématiques} publient-ils de la recherche, alors que leur avantage repose sur le secret~?
\end{interviewq}

\begin{interviewq}[$\star\star$ \iqroles{chercheur}]\label{iq:in:quantitative-hedge-funds-and-systematic-managers:5}
Les dix plus grands des 2\,577 conseillers détiennent un cinquième des actifs. Le secteur est-il concentré~? Qu'est-ce qui
changerait votre réponse~?
\end{interviewq}

\begin{interviewq}[$\star\star\star$ \iqroles{chercheur, risque}]\label{iq:in:quantitative-hedge-funds-and-systematic-managers:6}
Un programme de suivi de tendance et un livre d'arbitrage statistique ont chacun un ratio de Sharpe de un. Lequel peut recevoir
davantage de capital, et qu'est-ce qui limite chacun d'eux~?
\end{interviewq}
